\documentclass[11pt,a4paper]{article}
\usepackage[utf8]{inputenc}
\usepackage[vietnamese]{babel}
\usepackage{amsmath,amssymb}
\usepackage{geometry}
\geometry{a4paper, total={180mm,260mm}, left=15mm, right=15mm, top=18mm, bottom=18mm}
\usepackage{tikz}
\usetikzlibrary{shapes.geometric, arrows.meta, positioning, calc}
\usepackage{tcolorbox}
\tcbuselibrary{skins,breakable}
\usepackage{xcolor}
\usepackage{multicol}
\usepackage{array}
\usepackage{fontawesome}
\usepackage{fancyhdr}

\pagestyle{fancy}
\fancyhf{}
\renewcommand{\headrulewidth}{0.5pt}
\renewcommand{\footrulewidth}{0.5pt}
\fancyhead[L]{\small\bfseries\color{teal}\faBook\ SỔ TAY KIẾN THỨC TOÁN 3}
\fancyhead[R]{\small\bfseries\color{teal}Lớp Toán Thầy Lĩnh -- loptoanthaylinh.vn}
\fancyfoot[C]{\small\color{gray}-- Sổ tay kiến thức: Bài 77. Ôn tập phép cộng, phép trừ trong phạm vi 100 000 --}

\definecolor{primary}{HTML}{005F73}
\definecolor{secondary}{HTML}{0A9396}
\definecolor{accent}{HTML}{EE9B00}
\definecolor{danger}{HTML}{AE2012}
\definecolor{boxbg}{HTML}{F4F6F8}

\begin{document}

% --- BANNER TIÊU ĐỀ ---
\begin{tcolorbox}[
  enhanced,
  colback=white,
  colframe=primary,
  arc=3mm,
  boxrule=1.2pt,
  left=4mm, right=4mm, top=2.5mm, bottom=2.5mm,
  drop shadow=black!10!white
]
  \centering
  {\Large\bfseries\color{primary}\faBookmark\ BÀI 77: ÔN TẬP CỘNG TRỪ TRONG PHẠM VI 100 000}\\[1.5mm]
  {\small\itshape\color{black!75}Ôn tập cuối năm: Đặt tính thẳng cột, cộng trừ có nhớ, tính giá trị biểu thức và bài toán tìm thành phần chưa biết trên 1 trang A4}
\end{tcolorbox}

\vspace{-2mm}


\noindent
\begin{minipage}[t]{0.49\textwidth}\vspace{0pt}
\begin{tcolorbox}[
  enhanced,
  colback=blue!4!white,
  colframe=blue!75!black,
  arc=2mm, boxrule=0.8pt,
  left=2mm, right=2mm, top=1.8mm, bottom=1.8mm,
  title={\footnotesize\bfseries\color{white}\faCalculator\ 1. Quy tắc cộng và trừ}
]
{\footnotesize
\begin{itemize}\setlength{\itemsep}{1.5pt}\setlength{\leftskip}{-2.5mm}
  \item Đặt tính thẳng cột giữa các hàng từ phải sang trái.
  \item Thực hiện tính lần lượt từ \textbf{phải sang trái}:
  $$\text{ĐV} \to \text{Chục} \to \text{Trăm} \to \text{Nghìn} \to \text{Chục nghìn}$$
  \item \textbf{Cộng có nhớ:} Nhớ $1$ sang hàng liền kề bên trái.
  \item \textbf{Trừ có nhớ:} Mượn $1$ chục ở hàng bên trái, trả $1$ sang hàng tiếp theo của số trừ.
\end{itemize}
}
\end{tcolorbox}
\end{minipage}
\hfill
\begin{minipage}[t]{0.49\textwidth}\vspace{0pt}
\begin{tcolorbox}[
  enhanced,
  colback=teal!4!white,
  colframe=teal!75!black,
  arc=2mm, boxrule=0.8pt,
  left=2mm, right=2mm, top=1.8mm, bottom=1.8mm,
  title={\footnotesize\bfseries\color{white}\faSliders\ 2. Kĩ năng giải toán tổng hợp}
]
{\footnotesize
\begin{itemize}\setlength{\itemsep}{1.5pt}\setlength{\leftskip}{-2.5mm}
  \item \textbf{Tìm thành phần chưa biết:}\\
  $x + a = b \implies x = b - a$\\
  $x - a = b \implies x = b + a$\\
  $a - x = b \implies x = a - b$
  \item \textbf{Biểu thức có ngoặc:} Luôn tính trong ngoặc trước.
  \item Thử lại phép trừ bằng phép cộng tương ứng.
\end{itemize}
}
\end{tcolorbox}
\end{minipage}

\vspace{-2.5mm}
% --- KHỐI 3: SƠ ĐỒ TƯ DUY TIKZ ---
\noindent
\begin{tcolorbox}[
  enhanced,
  colback=orange!4!white,
  colframe=orange!80!black,
  arc=2mm, boxrule=0.8pt,
  left=2mm, right=2mm, top=1.5mm, bottom=1.5mm,
  title={\footnotesize\bfseries\color{white}\faSitemap\ SƠ ĐỒ TƯ DUY: HỆ THỐNG CỘNG TRỪ CUỐI NĂM LỚP 3}
]
\centering
\begin{tikzpicture}[
  box/.style={draw, rounded corners=2pt, align=center, font=\scriptsize\bfseries, inner sep=2.5pt, minimum height=0.75cm},
  arr/.style={-{Stealth[length=2mm]}, thick}
]
  \node[box, fill=purple!15, draw=purple!80!black, text width=2.6cm] (c1) at (0, 0) {ĐẶT TÍNH\\Thẳng cột tuyệt đối\\Tính từ phải qua trái};
  \node[box, fill=teal!15, draw=teal!80!black, text width=2.6cm] (c2) at (3.5, 0) {TÍNH NHẨM\\Số tròn chục nghìn\\Số tròn trăm nghìn};
  \node[box, fill=blue!12, draw=blue!80!black, text width=2.6cm] (c3) at (7.0, 0) {BIỂU THỨC SỐ\\Ngoặc tròn trước\\Tính từ trái qua phải};
  \node[box, fill=orange!15, draw=orange!80!black, text width=2.6cm] (c4) at (10.5, 0) {BÀI TOÁN 2 BƯỚC\\Tìm ẩn trung gian\\Tính tổng thể cuối};

  \draw[arr, purple!80!black] (c1) -- (c2);
  \draw[arr, teal!80!black] (c2) -- (c3);
  \draw[arr, blue!80!black] (c3) -- (c4);
\end{tikzpicture}
\end{tcolorbox}

\vspace{-2.5mm}
% --- KHỐI 4: VÍ DỤ & MẸO TRÁNH BẪY ---
\noindent
\begin{minipage}[t]{0.49\textwidth}\vspace{0pt}
\begin{tcolorbox}[
  enhanced,
  colback=green!4!white,
  colframe=green!60!black,
  arc=2mm, boxrule=0.8pt,
  left=2mm, right=2mm, top=1.8mm, bottom=1.8mm,
  title={\footnotesize\bfseries\color{white}\faCheckCircle\ Ví dụ thực hành chuẩn mực}
]
{\footnotesize
\textbf{1. Tìm $x$:}\\
$x + 35\,200 = 70\,000 \implies x = 70\,000 - 35\,200 = \mathbf{34\,800}$.\\
$82\,000 - x = 25\,000 \implies x = 82\,000 - 25\,000 = \mathbf{57\,000}$.\\
\textbf{2. Tính nhanh:}\\
$15\,000 + 42\,000 + 25\,000 = (15\,000 + 25\,000) + 42\,000$\\
$= 40\,000 + 42\,000 = \mathbf{82\,000}$.
}
\end{tcolorbox}
\end{minipage}
\hfill
\begin{minipage}[t]{0.49\textwidth}\vspace{0pt}
\begin{tcolorbox}[
  enhanced,
  colback=red!4!white,
  colframe=red!70!black,
  arc=2mm, boxrule=0.8pt,
  left=2mm, right=2mm, top=1.8mm, bottom=1.8mm,
  title={\footnotesize\bfseries\color{white}\faExclamationTriangle\ Mẹo tránh bẫy trắc nghiệm}
]
{\footnotesize
\begin{itemize}\setlength{\itemsep}{1.5pt}\setlength{\leftskip}{-2.5mm}
  \item \textbf{Bẫy mượn quên trả khi trừ có 2 lần nhớ liên tiếp:}\\
  Nếu mượn ở hàng đơn vị rồi tiếp tục mượn ở hàng chục, phải trả nhớ đầy đủ vào cả hai hàng tiếp theo của số trừ!
  \item Luôn kiểm tra lại số chữ số của hiệu và tổng.
\end{itemize}
}
\end{tcolorbox}
\end{minipage}


\end{document}
